% Build: make (latexmk in the TeX Live container; README.md)

\documentclass[journal]{IEEEtran}

\usepackage{cite}
\usepackage{amsmath,amssymb}
\usepackage{graphicx}
\usepackage{booktabs}
\usepackage{siunitx}
\usepackage{xcolor}
\usepackage[hidelinks]{hyperref}

% a result not measured yet (to be removed before submission)
\newcommand{\tbd}[1]{{\color{red!75!black}[#1]}}

\begin{document}

\title{Real-Time Transformer Digital Predistortion on an Embedded GPU for a 28-GHz Analog
Radio-over-Fiber Fronthaul}

\author{Ahmed~Abdellatif}

\markboth{Journal of Lightwave Technology,~Vol.~XX, No.~X, Month~2027}%
{Abdellatif: Real-Time Transformer DPD on an Embedded GPU for a 28-GHz A-RoF Fronthaul}

\maketitle

\begin{abstract}
\tbd{To be written last, from the results.}
\end{abstract}

\begin{IEEEkeywords}
Analog radio-over-fiber, millimeter wave, digital predistortion, transformer, indirect learning,
GPU, real time, fronthaul.
\end{IEEEkeywords}

\section{Introduction}\label{sec:intro}

\IEEEPARstart{A}{nalog} radio-over-fiber (A-RoF) carries the radio waveform itself on an optical
carrier from the central office (CO) to a remote antenna unit (RAU). The RAU stays simple and the
signal processing stays at the CO, which makes A-RoF a candidate fronthaul for the
\SI{28}{GHz} band of 5G. Because the waveform stays
analog, the distortions of the components accumulate along the link: the modulators at the CO, the
photodiode, and above all the millimeter-wave (mmW) power amplifier (PA) at the RAU, which
compresses the large amplitudes of a wideband signal and adds memory. On a \SI{28}{GHz} wideband
A-RoF fronthaul, compensating this nonlinearity with a neural network at the receiver improves on a
Volterra equalizer, at a higher complexity~\cite{song2024}; but a receiver-side compensator has to
run in every user equipment, after the wireless hop has added its noise to the distortion.

Digital predistortion (DPD) compensates the same distortion once, at the CO, before it occurs: the
baseband drive is distorted with the inverse of the link, so that the PA's output is linear.
Indirect learning makes this practical: a postdistorter is trained to map the link's output back to
its input and is then copied in front of the link. With memory and generalized memory polynomials
identified in this way, DPD has linearized A-RoF links~\cite{hadi2019}. Polynomial models, however,
grow quickly with the order of the nonlinearity and the length of the memory, and neural networks
model strong nonlinearities with memory more compactly, at the cost of more computation per sample,
which a DPD must perform at the signal's rate. A GPU can carry that cost: a pruned feed-forward DPD
network has run in real time on a GPU at about \SI{1}{GSa/s}~\cite{li2025}. A GPU is also what a CO
or an edge node is likely to have, and an embedded GPU fits beside the fronthaul equipment.

This paper presents a transformer DPD for a \SI{28}{GHz} A-RoF fronthaul, running in real time on
an embedded GPU inside a real-time testbed of the link. The link uses a single laser: an optical
double-sideband signal generates the mmW carrier on the photodiode, and 16QAM at
\SI{400}{MBd} is modulated at baseband onto the upper sideband only, over \SI{15}{km} of
standard fiber, to a PA at the RAU. A Raspberry Pi emulates the link sample by sample as a pipeline
of stages; the DPD replaces the pass-through stage before the modulators and runs on a Jetson Orin
Nano. The transformer, a window of 32 samples encoded by two attention layers with
\num{18338} parameters, is trained offline by indirect learning on recordings of the live
testbed, and runs as hand-written CUDA kernels in a C++ stage of the pipeline.

The contributions of this paper are:
\begin{enumerate}
\item \emph{A transformer DPD for a mmW A-RoF link}, trained by indirect learning on the link's
own recordings, against a memory-polynomial DPD and no DPD on the same link
(Sections~\tbd{}): \tbd{the EVM and BER at the receiver, and the PA's output}.
\item \emph{Its real-time implementation on an embedded GPU}: the model as hand-written CUDA
kernels, checked against the PyTorch reference, with \tbd{its throughput and latency per frame}
(Section~\tbd{}).
\item \emph{A real-time testbed} in which the DPD runs as one stage of the emulated link, with the
link's other stages on a second machine and the two clocks synchronized, so that the DPD is trained
and evaluated on the stream it predistorts (Section~\tbd{}).
\end{enumerate}

Section~\tbd{} describes the link and the testbed, Section~\tbd{} the DPD and its training,
Section~\tbd{} its implementation on the GPU, Section~\tbd{} the results, and
Section~\tbd{} concludes.

\bibliographystyle{IEEEtran}
\bibliography{refs}

\end{document}
